% ═════════════════════════════════════════════════════════════════════════════
\begin{document}
\twocolumn[
\begin{@twocolumnfalse}
\begin{center}
{\LARGE\bfseries\color{aksen}
        Kecerdasan Buatan dalam Pendidikan Modern:
        Peluang, Tantangan, dan Arah Pengembangan
      }

\vspace{10pt}
      {\large Ahmad Fauzi$^{1}$, Siti Rahayu$^{2}$, Budi Santoso$^{3}$}
  
      \vspace{6pt}
      {$^{1}$Departemen Ilmu Komputer, Universitas Indonesia, Jakarta\newline 
        $^{2}$Fakultas Keguruan dan Ilmu Pendidikan, Universitas Gadjah Mada, Yogyakarta\newline 
        $^{3}$Program Studi Teknologi Pendidikan, Institut Teknologi Bandung, Bandung

\vspace{4pt}
        \textit{Korespondensi: \href{mailto:ahmad.fauzi@ui.ac.id}{ahmad.fauzi@ui.ac.id}}
      }

\vspace{10pt}
      \hrule height 1pt
      \vspace{8pt}
      \begin{abstract}
Artikel ini mengkaji peran kecerdasan buatan (AI)...
\end{abstract}
      \vspace{8pt}
      \hrule height 0.5pt
      \vspace{12pt}
\end{center}
\end{@twocolumnfalse}
]